
\frontslide{Structures d'indexation}


\begin{frame}[fragile]{Structures d'index}

Contenu de ce cours\,:

\begin{redbullet}
   \item Qu'est-ce qu'un index?
   
   \item Principes généraux
   
   \item Index denses et non-denses
\end{redbullet}

\vfill

\alert{En l'absence d'un index}, seule solution\,: le parcours séquentiel du fichier.
\vfill

\alert{Avec un index}: \alert{accès direct à l'enregistrement}, amélioration \alert{très importante} des temps de réponse


\vfill


\begin{blocgauche}
\alert{Ces diapositives correspondent au support en ligne disponible sur
le site \url{http://sys.bdpedia.fr/arbreb.html\#s1-indexation-de-fichiers}}
\end{blocgauche}

\end{frame}


%------------------------------------------------------------
\begin{frame}{Exemple 1\,: une table}

\begin{center}
\begin{small}
\begin{tabular}{l|c|c||l|c|c}
titre & année  & ...  & titre & année  & ...  \\
\hline
Vertigo  & 1958 & ... & Annie Hall & 1977 & ... \\
Brazil & 1984 & ... & Jurasic Park & 1992 & ... \\
Twin Peaks & 1990 & ...&  Metropolis & 1926 & ... \\
Underground & 1995 & ... & Manhattan & 1979 & ... \\
Easy Rider & 1969 & ... & Reservoir Dogs & 1992 & ... \\
Psychose & 1960 & ... & Impitoyable & 1992 & ... \\
Greystoke & 1984 & ... & Casablanca & 1942 & ... \\
Shining & 1980 & ... & Smoke & 1995 & ... \\
\hline
\end{tabular}
\end{small}
\end{center}

\end{frame}
%------------------------------------------------------------
\begin{frame}{Exemple 2}

La table des films,  avec\,:
\begin{redbullet}
  \item 1\,000\,000 (un million) de films
  \item Un enregistrement = 1200 octets
  \item Un bloc = 4K =$>$ 3 enregistrements par bloc 
  \item environ 300\,000 blocs, 1,2 Go
\end{redbullet}

\end{frame}
%------------------------------------------------------------
\section{Principes}
%------------------------------------------------------------


\begin{frame}{Index = concept bien connu}

Prenons un livre à contenu technique (pas une fiction!), par
exemple un livre de recettes. Il contient (au moins) un index.

\vfill
\begin{blocgauche}
\begin{redbullet}
  \item L'index présente les termes importants, classés par ordre alphabétique
  
  \item À chaque terme sont associés les numéros de page où on trouve le terme.
  
  \item En parcourant l'index (par dichotomie!) on trouve la ou les pages qui nous intéressent.
\end{redbullet}
\end{blocgauche}
\end{frame}

\begin{frame}{Les index informatiques}

Même principe! C'est un fichier qui permet de trouver 
un enregistrement dans une table.  
 Vocabulaire:

\vfill
\begin{blocgauche}
\begin{redbullet}
  \item \alert{Clé d'indexation} = une \alert{liste}
     d'un ou plusieurs attributs.
   \item Une \alert{adresse} (déjà vu) est une adresse de bloc ou une adresse d'enregistrement.
   \item \alert{Entrée d'index}: enregistrements de la forme \texttt{[valeur, addr]}, \texttt{valeur}
     est une valeur de clé. 
  \item L'index est \alert{trié} sur \texttt{valeur}

\end{redbullet}
\end{blocgauche}

\end{frame}
%------------------------------------------------------------
\begin{frame}{Exemples}

Clés de recherche\,:
\begin{redbullet}
  \item Le titre du film (c'est aussi la clé primaire)
  \item l'année du film 
  \item Une paire constituée du  titre et de l'année  
\end{redbullet}

\vfill

Opérations\,:
\begin{redbullet}
  \item Rechercher \textit{Vertigo}
  \item Rechercher les films parus entre 1960 et 1975
  \item Rechercher les films commençant par 'V'
\end{redbullet}

\vfill
\begin{blocgauche}
\alert{L'index ne sert à rien pour toute recherche
ne portant pas sur la clé}.
\end{blocgauche}
\end{frame}

\begin{frame}[fragile]{Le cas des dictionnaires}

Les dictionnaires ont la particularité de \alert{trier} les termes.

\vfill

On peut alors créer un index qui ne référence que \alert{le premier mot de chaque page}.

\begin{minted}{text}
  - ...
  - ballon, page 56
  - bille, page 57
  - bulle, page 65,
  - cable, page 72
  - ...
\end{minted}

\begin{blocgauche}
\alert{On peut quand même utiliser cet index pour trouver n'importe quel mot}.
("armée", "crabe", "botte", "belle").
\end{blocgauche}
\end{frame}
%%%%%%%%%
\section{Index non dense}
%------------------------------------------------------------

%------------------------------------------------------------
\begin{frame}{Index non dense}


 \figSlideWithSize{indexNonDense}{13cm}

\vfill
\begin{blocgauche}
\alert{Hypothèse}\,: le fichier de données est \alert{trié sur la clé},
comme un dictionnaire. \alert{L'index ne référence que la première
valeur de chaque bloc}.
\end{blocgauche}
\end{frame}
%------------------------------------------------------------
\begin{frame}{Opérations}

Recherches\,:
\begin{redbullet}
  \item \textbf{Par clé}\,: je recherche \textit{Shining}
  \item \textbf{Par intervalle}\,: tous les films entre \textit{Greystocke}
    et \textit{Psychose}.

  \item \textbf{Par préfixe}\,: tous les films commençant par 'M'.
\end{redbullet}

\vfill
\begin{blocgauche}
\alert{Mais pas par suffixe}: tous les films terminant par 'e'?
\end{blocgauche}
\end{frame}
%------------------------------------------------------------
\begin{frame}{Exemple concret}

Sur notre fichier de 1,2 Go
\begin{redbullet}
  \item En supposant qu'un titre occupe 20 octets, une adresse 8 octets
   \item Taille de l'index\,: $300\,000 * (20 + 8)=8,4 Mo$ octets
\end{redbullet}

\vfill
\begin{blocgauche}
Beaucoup plus petit que le fichier\,!

\medskip

Problème\,: maintenir l'ordre sur le fichier \alert{et}
sur l'index.
\end{blocgauche}

\end{frame}
\section{Index dense}
%------------------------------------------------------------
\begin{frame}{Index dense}

  \figSlideWithSize{indexDense}{13cm}

\vfill

\begin{blocgauche}
\alert{Fichier de données non trié. Toutes les valeurs de clé sont représentées}
\end{blocgauche}

\end{frame}
%------------------------------------------------------------
\begin{frame}{Exemple concret}

Sur notre fichier de 1,2 Go
\begin{redbullet}
  \item Une année = 4 octets, une adresse 8 octets
   \item Taille de l'index\,: $1\,000\,000 * (4 + 8)=12$ Mo
\end{redbullet}

Encore 100 fois plus petit que le fichier.

\vfill

Recherches\,:
\begin{redbullet}
  \item Par clé\,: comme sur un index non-dense

   \item Par intervalle (exemple [1950, 1979])\,:
     \begin{itemize}
       \item recherche, dans l'index de la borne inférieure
       \item parcours séquentiel {\em dans l'index}
      \item à chaque valeur\,: accès au fichier de données
   \end{itemize}
   \alert{Engendre des accès aléatoires.}
\end{redbullet}


\end{frame}
\subsection{Index multi-niveaux}
%------------------------------------------------------------
\begin{frame}{Index multi-niveaux: on indexe l'index}


  \figSlideWithSize{index2Niveaux}{12cm}

\vfill
\begin{blocgauche}
\begin{redbullet}
  \item \alert{Essentiel}\,: l'index est trié,
     donc on peut l'indexer par un second niveau
          \alert{non-dense}
   \item Sinon ça ne servirait à rien (pourquoi\,?)
 \end{redbullet}
\end{blocgauche}

\end{frame}
%------------------------------------------------------------
\begin{frame}{Index multi-niveaux: jusqu'où?}


  \figSlideWithSize{indexMultiNiv}{12cm}

\vfill
\begin{blocgauche}
\begin{redbullet}
  \item Arrêt quand racine constitué d'un seul bloc.
   \item structure \alert{hiérarchique}; recherche
    \alert{de haut en bas}.
 \end{redbullet}
\end{blocgauche}

\end{frame}


\begin{frame}{Résumé}

\vfill 
\alert{Retenir}:

\begin{redbullet}
  \item Un index accélère des opérations de recherche portant sur la \alert{clé d'indexation} ou sur un préfixe de la clé
  \item La structure repose sur le tri des valeurs de clé indexées
  \item On peut créer un \alert{index non dense} sur un fichier trié, \alert{dense} sur un fichier non-trié.
  \item Les index multi-niveaux s'appuient sur un premier niveau dense, puis sur des niveaux supérieurs non-denses.
\end{redbullet}

\vfill

\begin{blocgauche}
\alert{Les index ont un coût}, il faut les maintenir au fur et à mesure des insertions dans le fichier indexé.
\end{blocgauche}

\end{frame}
